\documentclass[a4paper,12pt]{article}

\usepackage[utf8]{inputenc}
\usepackage[T1]{fontenc}
\usepackage[italian]{babel}
\usepackage{amsmath}
\usepackage{accsupp}
\usepackage{hyperref}

% Comando per associare testo alternativo/accessibile a una formula
\newcommand{\altmath}[2]{%
  \BeginAccSupp{method=pdfstringdef,ActualText={#2}}%
  #1%
  \EndAccSupp{}%
}

\title{Esempio di formula matematica accessibile in LaTeX}
\author{}
\date{}

\begin{document}

\maketitle

Questo documento mostra un esempio minimale di formula matematica resa più accessibile
tramite testo alternativo nel PDF.

La formula seguente è visualizzata normalmente, ma contiene anche un testo alternativo
leggibile dagli strumenti di accessibilità:

\[
\altmath{x^2 + 3x - 4 = 0}{x al quadrato più 3x meno 4 uguale zero}
\]

Questo permette a un lettore PDF compatibile di esporre una descrizione testuale
più comprensibile della formula.

\end{document}